\scanpage{16}
\mspage{\handwritten{15}}

\deleted{tés, les morphismes \uline{toutes} les correspondances \inserted{co}homologiques (en \unsure{vertu} du thorème de Thom rappelé plus haut). On trouve donc un foncteur pleinement fidèle de cette catégorie dans celle des vectoriels gradués \deleted{\unsure{de d|degr}} positivement de dim. finie sur $\underline F_2$, qui est déjà abélienne, et si on veut rendre comme plus haut $\underline F_2(-1)$ inversible, on \unsure{trouve} tous les vectoriels gradués de dim finie sur $\underline F_2$. C'est un peu \unsure{trop} trivial comme structure, et on peut faire mieux,} \deleted{en prenant comme correspondances \unsure{pour ceci|Pour ceci}}

\inserted{Pour trouver une catégorie universelle,} on associe à toute $X$ \deleted{le sous-anneau \unsure{ensemble} de} l'ensemble \inserted{$S(X)$} des sous-variétés \deleted{$S(X)$} de $X$, modulo équivalence homologique. Donc $S(X)\handwritten{\subset}H^{\bullet}(X)$, et on prouve qu'on trouve en fait un sous-anneau de $H^{\bullet}(X)$ (grâce à un théorème de Thom qui caractérise $S(X)$ \handwritten{en} termes d'opérations de Steenrod, cf plus bas). On prouve de même (par le même théorème) que les $S(X\mathbin{x}Y)$ sont stables par composition \unsure{de} correspondances cohomologiques\handwritten{(} cela ne semble pas trivial, puisque les $S(X)$ ne sont pas stables par Gysin...). D'où une catégorie additive dont les objets sont ceux de $C$, les flêches sont les \inserted{\handwritten{él. des}} $S(X\mathbin{x}Y)$, Il semble qu'elle doive jouer le rôle universel que j'ai envisagé plus haut. En tous cas, sa définition géo\deleted{\unsure{é}}métrique est fort simple Nous allons en donner une interpétation algébrique. Pŏur ce\unsure{ĉ|c}\unsure{î|i}, considérons l'anneau des opérations cohomologiques \deleted{additives sur \unsure{i.}} \inserted{``stables''} \deleted{des endomorphismes du \deleted{\unsure{le|la}} foncteur $X\longrightarrow H^{\bullet}(X)$ qui sont additifs, et qui en degré zéro se réduisent à un multiple de l'identité.} C'est un anneau gradué, avec application diagonale anticommutative\inserted{ i.e. commutative (car. $2$ !)}, associative et unitaire. Son dual gradué est une algèbre \inserted{graduée \unsure{antioommutative|anticommutative} donc \deleted{commutative},} commutative,avec application diagonale associative et unitaire; il est réduit à $\underline F_2$ en degré zéro. Donc il corespond à un schéma en groupes affine $G$ sur $\underline F_2$, qui est \unsure{d'ailleurs} le foncteur des $\otimes$-automorphismes \inserted{additifs} du foncteur $H^{\bullet}$ sur $C$ \inserted{\handwritten{qui \uline{opèrent} \unsure{trivialement} sur \deleted{\unsure{l'élément}} $\unsure{H^1(S^1)|H^1(\Delta^1)}$. \deleted{\unsure{On}}}}
\margin{\handwritten{\unsure{?|??} \(\vert\)}}% at the left of the definition of the stable operations

% archive sheet number: 15
